\documentclass[reqno]{amsart}
\usepackage[T1]{fontenc}
\usepackage{amsmath,amssymb,amsthm}
\usepackage{xr-hyper}
\usepackage[hidelinks]{hyperref}
\externaldocument[da-]{20261126T030300--the-data__workbook}
\externaldocument[ob-]{20261126T040100--spaces-forms-and-twists__workbook}
\externaldocument[cs-]{20261125T003000--the-chain-and-the-square-root__workbook}
\externaldocument[dch-]{20261126T030700--the-determinant-of-a-chain__workbook}
\externaldocument[c1-]{20261126T050400--the-norm-of-one-lift__workbook}
\externaldocument[c2-]{20261126T050500--equal-multiplicities__workbook}
\externaldocument[c3-]{20261126T050600--the-exchange-identity__workbook}
\externaldocument[tc-]{20261126T040200--theorem-c__workbook}

\newcounter{chap}
% Chapter numbers, generated from 20261125T012000--grh-workbook__guide.org
% by scripts/make-chapters.py.  Do not edit: change the order of the
% index and rerun the script (./preview-site.sh and CI do this).
%
% Usage in a sheet, after \newcounter{chap}:
%   % Chapter numbers, generated from 20261125T012000--grh-workbook__guide.org
% by scripts/make-chapters.py.  Do not edit: change the order of the
% index and rerun the script (./preview-site.sh and CI do this).
%
% Usage in a sheet, after \newcounter{chap}:
%   % Chapter numbers, generated from 20261125T012000--grh-workbook__guide.org
% by scripts/make-chapters.py.  Do not edit: change the order of the
% index and rerun the script (./preview-site.sh and CI do this).
%
% Usage in a sheet, after \newcounter{chap}:
%   \input{chapters}\chapnumber{<sheet ID>}
\makeatletter
\providecommand\chapters@set[2]{\@namedef{chapters@#1}{#2}}
\chapters@set{20261125T001000}{1}% zeros-to-ramp-sums
\chapters@set{20261125T001100}{2}% inert-primes-and-estimate-b
\chapters@set{20261125T005000}{3}% galois-fourier-and-averaging
\chapters@set{20261125T005500}{4}% columns-from-a-sphere
\chapters@set{20261125T000000}{5}% slope-method
\chapters@set{20261126T030100}{6}% a-twist-for-each-prime
\chapters@set{20261125T003000}{7}% the-chain-and-the-square-root
\chapters@set{20261125T010000}{8}% rank-lemma-and-zero-estimate
\chapters@set{20261125T010500}{9}% poincare-series-and-near-orthogonality
\chapters@set{20261124T223000}{10}% elliptic-curves-for-the-proof
\chapters@set{20261125T009500}{11}% fock-model-and-theta-intertwiners
\chapters@set{20261124T223100}{12}% complex-multiplication-and-inert-reduction
\chapters@set{20261125T008000}{13}% one-cm-determinant
\chapters@set{20261124T214500}{14}% supersingular-theta-divisibility
\chapters@set{20261124T214600}{15}% the-isotropic-plane
\chapters@set{20261125T013700}{16}% pairings-on-a-formal-group-of-height-two
\chapters@set{20261125T013500}{17}% hecke-nondegeneracy-on-the-tree
\chapters@set{20261126T030200}{18}% the-bound-is-attained
\chapters@set{20261125T009000}{19}% kronecker-limit-and-chowla-selberg
\chapters@set{20261125T002000}{20}% definite-theta-lifts
\chapters@set{20261125T015500}{21}% siegel-weil-in-the-smallest-case
\chapters@set{20261125T002100}{22}% definite-theta-lifts-and-norms
\chapters@set{20261125T016500}{23}% hecke-operators-and-theta-lifts
\chapters@set{20261125T016700}{24}% divisibility-along-the-chain
\chapters@set{20261125T004000}{25}% denominators-and-sha
\chapters@set{20261125T004100}{26}% tate-twists-and-sha
\chapters@set{20261126T020430}{27}% the-minus-part-by-gauss-sums
\chapters@set{20261125T014100}{28}% congruence-modules
\chapters@set{20261125T014500}{29}% eisenstein-congruences-and-selmer-classes
\chapters@set{20261125T014700}{30}% pseudocharacters-and-the-length-bound
\chapters@set{20261125T011000}{31}% galois-representations-black-box
\chapters@set{20261126T030300}{32}% the-data
\chapters@set{20261126T020210}{33}% bessel-poincare-series-and-their-gram-operators
\chapters@set{20261126T030400}{34}% the-seesaw-identity
\chapters@set{20261126T030500}{35}% divisibility-of-the-composite-lift
\chapters@set{20261126T030600}{36}% theorem-a-short-forms
\chapters@set{20261126T030700}{37}% the-determinant-of-a-chain
\chapters@set{20261126T030800}{38}% the-cost-of-one-lift
\chapters@set{20261126T030900}{39}% denominators-of-projections-to-lifts
\chapters@set{20261126T031000}{40}% twists-with-few-congruences
\chapters@set{20261126T031100}{41}% theorem-c-and-the-main-inequality
\providecommand\chapnumber[1]{%
  \@ifundefined{chapters@#1}%
    {\PackageError{chapters}{Sheet #1 is not listed in the index}%
      {List it in 20261125T012000--grh-workbook__guide.org\MessageBreak
       and run scripts/make-chapters.py.}}%
    {\edef\thechap{\@nameuse{chapters@#1}}}}
\makeatother
\chapnumber{<sheet ID>}
\makeatletter
\providecommand\chapters@set[2]{\@namedef{chapters@#1}{#2}}
\chapters@set{20261125T001000}{1}% zeros-to-ramp-sums
\chapters@set{20261125T001100}{2}% inert-primes-and-estimate-b
\chapters@set{20261125T005000}{3}% galois-fourier-and-averaging
\chapters@set{20261125T005500}{4}% columns-from-a-sphere
\chapters@set{20261125T000000}{5}% slope-method
\chapters@set{20261126T030100}{6}% a-twist-for-each-prime
\chapters@set{20261125T003000}{7}% the-chain-and-the-square-root
\chapters@set{20261125T010000}{8}% rank-lemma-and-zero-estimate
\chapters@set{20261125T010500}{9}% poincare-series-and-near-orthogonality
\chapters@set{20261124T223000}{10}% elliptic-curves-for-the-proof
\chapters@set{20261125T009500}{11}% fock-model-and-theta-intertwiners
\chapters@set{20261124T223100}{12}% complex-multiplication-and-inert-reduction
\chapters@set{20261125T008000}{13}% one-cm-determinant
\chapters@set{20261124T214500}{14}% supersingular-theta-divisibility
\chapters@set{20261124T214600}{15}% the-isotropic-plane
\chapters@set{20261125T013700}{16}% pairings-on-a-formal-group-of-height-two
\chapters@set{20261125T013500}{17}% hecke-nondegeneracy-on-the-tree
\chapters@set{20261126T030200}{18}% the-bound-is-attained
\chapters@set{20261125T009000}{19}% kronecker-limit-and-chowla-selberg
\chapters@set{20261125T002000}{20}% definite-theta-lifts
\chapters@set{20261125T015500}{21}% siegel-weil-in-the-smallest-case
\chapters@set{20261125T002100}{22}% definite-theta-lifts-and-norms
\chapters@set{20261125T016500}{23}% hecke-operators-and-theta-lifts
\chapters@set{20261125T016700}{24}% divisibility-along-the-chain
\chapters@set{20261125T004000}{25}% denominators-and-sha
\chapters@set{20261125T004100}{26}% tate-twists-and-sha
\chapters@set{20261126T020430}{27}% the-minus-part-by-gauss-sums
\chapters@set{20261125T014100}{28}% congruence-modules
\chapters@set{20261125T014500}{29}% eisenstein-congruences-and-selmer-classes
\chapters@set{20261125T014700}{30}% pseudocharacters-and-the-length-bound
\chapters@set{20261125T011000}{31}% galois-representations-black-box
\chapters@set{20261126T030300}{32}% the-data
\chapters@set{20261126T020210}{33}% bessel-poincare-series-and-their-gram-operators
\chapters@set{20261126T030400}{34}% the-seesaw-identity
\chapters@set{20261126T030500}{35}% divisibility-of-the-composite-lift
\chapters@set{20261126T030600}{36}% theorem-a-short-forms
\chapters@set{20261126T030700}{37}% the-determinant-of-a-chain
\chapters@set{20261126T030800}{38}% the-cost-of-one-lift
\chapters@set{20261126T030900}{39}% denominators-of-projections-to-lifts
\chapters@set{20261126T031000}{40}% twists-with-few-congruences
\chapters@set{20261126T031100}{41}% theorem-c-and-the-main-inequality
\providecommand\chapnumber[1]{%
  \@ifundefined{chapters@#1}%
    {\PackageError{chapters}{Sheet #1 is not listed in the index}%
      {List it in 20261125T012000--grh-workbook__guide.org\MessageBreak
       and run scripts/make-chapters.py.}}%
    {\edef\thechap{\@nameuse{chapters@#1}}}}
\makeatother
\chapnumber{<sheet ID>}
\makeatletter
\providecommand\chapters@set[2]{\@namedef{chapters@#1}{#2}}
\chapters@set{20261125T001000}{1}% zeros-to-ramp-sums
\chapters@set{20261125T001100}{2}% inert-primes-and-estimate-b
\chapters@set{20261125T005000}{3}% galois-fourier-and-averaging
\chapters@set{20261125T005500}{4}% columns-from-a-sphere
\chapters@set{20261125T000000}{5}% slope-method
\chapters@set{20261126T030100}{6}% a-twist-for-each-prime
\chapters@set{20261125T003000}{7}% the-chain-and-the-square-root
\chapters@set{20261125T010000}{8}% rank-lemma-and-zero-estimate
\chapters@set{20261125T010500}{9}% poincare-series-and-near-orthogonality
\chapters@set{20261124T223000}{10}% elliptic-curves-for-the-proof
\chapters@set{20261125T009500}{11}% fock-model-and-theta-intertwiners
\chapters@set{20261124T223100}{12}% complex-multiplication-and-inert-reduction
\chapters@set{20261125T008000}{13}% one-cm-determinant
\chapters@set{20261124T214500}{14}% supersingular-theta-divisibility
\chapters@set{20261124T214600}{15}% the-isotropic-plane
\chapters@set{20261125T013700}{16}% pairings-on-a-formal-group-of-height-two
\chapters@set{20261125T013500}{17}% hecke-nondegeneracy-on-the-tree
\chapters@set{20261126T030200}{18}% the-bound-is-attained
\chapters@set{20261125T009000}{19}% kronecker-limit-and-chowla-selberg
\chapters@set{20261125T002000}{20}% definite-theta-lifts
\chapters@set{20261125T015500}{21}% siegel-weil-in-the-smallest-case
\chapters@set{20261125T002100}{22}% definite-theta-lifts-and-norms
\chapters@set{20261125T016500}{23}% hecke-operators-and-theta-lifts
\chapters@set{20261125T016700}{24}% divisibility-along-the-chain
\chapters@set{20261125T004000}{25}% denominators-and-sha
\chapters@set{20261125T004100}{26}% tate-twists-and-sha
\chapters@set{20261126T020430}{27}% the-minus-part-by-gauss-sums
\chapters@set{20261125T014100}{28}% congruence-modules
\chapters@set{20261125T014500}{29}% eisenstein-congruences-and-selmer-classes
\chapters@set{20261125T014700}{30}% pseudocharacters-and-the-length-bound
\chapters@set{20261125T011000}{31}% galois-representations-black-box
\chapters@set{20261126T030300}{32}% the-data
\chapters@set{20261126T020210}{33}% bessel-poincare-series-and-their-gram-operators
\chapters@set{20261126T030400}{34}% the-seesaw-identity
\chapters@set{20261126T030500}{35}% divisibility-of-the-composite-lift
\chapters@set{20261126T030600}{36}% theorem-a-short-forms
\chapters@set{20261126T030700}{37}% the-determinant-of-a-chain
\chapters@set{20261126T030800}{38}% the-cost-of-one-lift
\chapters@set{20261126T030900}{39}% denominators-of-projections-to-lifts
\chapters@set{20261126T031000}{40}% twists-with-few-congruences
\chapters@set{20261126T031100}{41}% theorem-c-and-the-main-inequality
\providecommand\chapnumber[1]{%
  \@ifundefined{chapters@#1}%
    {\PackageError{chapters}{Sheet #1 is not listed in the index}%
      {List it in 20261125T012000--grh-workbook__guide.org\MessageBreak
       and run scripts/make-chapters.py.}}%
    {\edef\thechap{\@nameuse{chapters@#1}}}}
\makeatother
\chapnumber{20261126T050700}
\renewcommand{\thesection}{\thechap.\arabic{section}}

\theoremstyle{definition}
\newtheorem{problem}{Problem}[section]
\newtheorem{definition}{Definition}[section]
\theoremstyle{plain}
\newtheorem{theorem}{Theorem}[section]

\newcommand{\Z}{\mathbb Z}
\newcommand{\R}{\mathbb R}
\newcommand{\C}{\mathbb C}
\newcommand{\F}{\mathbb F}
\newcommand{\Nm}{\operatorname{Nm}}
\newcommand{\Rc}{\mathcal R}
\newcommand{\cM}{\mathcal M}
\newcommand{\Oc}{\mathcal O}
\newcommand{\GL}{\mathrm{GL}}
\newcommand{\ex}{\mathrm{ex}}
\newcommand{\cmat}{c_{\mathrm{mat}}}
\newcommand{\cwt}{c_{\mathrm{wt}}}
\newcommand{\Klast}{K_{\mathrm{last}}}
\newcommand{\Hc}{H_\circ}
\newcommand{\tl}[1]{\texttt{#1}}
\setlength{\emergencystretch}{3em}

\title{Squares and the triangle}
\author{}
\date{November 26, 2026}

\begin{document}
\maketitle
\refstepcounter{chap}\label{chap}

\paragraph{Prerequisites.}
Inner products are conjugate-linear in the first argument.  From Chapter~\ref{c1-chap}: Theorem~\ref{c1-thm:one-lift}, the twist factor \(c_Q\) (Definition~\ref{c1-def:twist-factor}), the Rallis operators \(\Rc_j(\xi; V)\) (Definition~\ref{c1-def:rallis-operator}) and Problem~\ref{c1-pr:left-inverse}.  From Chapter~\ref{c3-chap}: moduli \(\mathfrak m\) and \(S_{\mathfrak m}\) (Definition~\ref{c3-def:modulus}), the exchange factor \(\ex_h(S, \nu)\) (Definition~\ref{c3-def:exchange-factor}), Theorem~\ref{c3-thm:exchange} and Problem~\ref{c3-pr:ex-ratio}.  From Chapter~\ref{ob-chap}: \(\cM\), \(k_w\), \(q_w\), the spaces \(Y_j\) (rank \(r + jt\)), \(H(V) \supseteq H_g(V)\), \(T\), \(m\)-fold twists \(\vartheta\) with tame part \((\psi_w)\) and set \(Q_\vartheta\), the tame part \((\omega_w)\) of \(\chi_M\), of order \(q_2\), the characters \(\mu_j^0\), the lower chain \(\Theta_\vartheta\) (Section~\ref{ob-sec:chain}), and \(L^S(z, \nu)\) and critical pairs (Definition~\ref{ob-def:pair}).  Chapter~\ref{da-chap} supplies \(t = 5r\) and \(b\), and Section~\ref{ob-sec:weights} defines \(\Delta = 2r^8\).  Chapter~\ref{cs-chap} is motivation.

\medskip
Theorem~\ref{dch-thm:main} factors \(Z(\vartheta)/Z(1)\).  It uses the injectivity of the chain (Problem~\ref{pr:theta-injective}) and the Gram identity of Theorem~\ref{thm:chain-gram}.  We derive the Gram identity from the realization of a triangle of lifts, the inner product formula on its first edges, and the comparison at the last level, all three read from the proof.  The square laws in the realization of the triangle are what the proof obtains from the exchange identity (Theorem~\ref{c3-thm:exchange}, stated and not proved in Chapter~\ref{c3-chap}).  By Theorem~\ref{c1-thm:one-lift}, for each \(1 \le j \le m\) there is a nonzero \(a_j \in \C\) such that for every tame character \(\xi\) the Gram operator of the single lift on \(H(W_j)\) is \(a_j c_{Q_\xi}\Rc_j(\xi; W_j)\).  Why is the Gram operator of the chain on \(T\) a scalar multiple of the product of the \(m\) Rallis operators, and how do ratios of \(L\)-values enter that scalar?

\section{Why a triangle}\label{sec:why}

\begin{problem}\label{pr:range}
With \(t = 5r\): \(2r < r + t\), but \(2(r + (j-1)t) > r + jt\) for every \(j \ge 2\).  So the condition \(2n^\flat < n^\sharp\) of the inner product formula (Theorem~\ref{c1-thm:rallis}) holds for the first lift \(Y_0 \to Y_1\) of the chain and for no later lift \(Y_{j-1} \to Y_j\).
\end{problem}

The following identity would suffice if the Gram operators could be moved through the preceding lifts.

\begin{problem}\label{pr:push-down}
Let \(u_j \colon \mathcal V_{j-1} \to \mathcal V_j\) (\(1 \le j \le m\)) be linear maps of complex inner product spaces that have adjoints, and \(U_j := u_j \circ \dots \circ u_1\).  Suppose that \(u_1^*u_1 = a_1 \mathcal E_1\) and that \(u_j^*u_j\,U_{j-1} = a_j\,U_{j-1} \mathcal E_j\) for \(2 \le j \le m\), with scalars \(a_j\) and endomorphisms \(\mathcal E_j\) of \(\mathcal V_0\).  Then \(U_m^*U_m = a_1 \cdots a_m\, \mathcal E_1 \mathcal E_2 \cdots \mathcal E_m\).
\end{problem}

For \(2 \le j \le m\) the ranks of the chain are outside the range of the inner product formula (Problem~\ref{pr:range}).  The proof applies the formula only on the first edge of each row of a triangle of lifts, and relates the resulting Gram forms by two square laws.

\section{The triangle and its two square laws}\label{sec:triangle}

In this sheet a \emph{triangle of size \(m\)} consists of inner product spaces \(H_{i,n}\) (\(i, n \ge 0\), \(i + n \le m\)); a space \(\Hc\) with linear isometric embeddings \(\Hc \to H_{i,0}\), through which we regard \(\Hc\) as a subspace of each \(H_{i,0}\); linear maps \(\lambda_{i,n} \colon H_{i,n} \to H_{i,n+1}\) (lifts) and \(\delta_{i,n} \colon H_{i,n+1} \to H_{i+1,n}\) (descents) for \(i + n < m\); and positive integer weights \(g_0, g_1, \dots\).  The \emph{weighted rows} are \(\lambda_i^{(0)} := 1\) and \(\lambda_i^{(n+1)} := g_n \lambda_{i,n} \lambda_i^{(n)} \colon H_{i,0} \to H_{i,n+1}\).  The \emph{square} \((i,n)\), for \(i + n + 2 \le m\), has corners \(H_{i,n+1}\), \(H_{i,n+2}\), \(H_{i+1,n}\), \(H_{i+1,n+1}\).  We ask:
\begin{itemize}
\item[(F)] for \(i < m\), \(x \in \Hc\) and \(v \in H_{i,1}\): \(\langle v, \lambda_{i,0}x\rangle = \langle \delta_{i,0}v, x\rangle\); and \(\delta_{i,0}\lambda_{i,0}\) maps \(\Hc \subseteq H_{i,0}\) injectively into \(\Hc \subseteq H_{i+1,0}\);
\item[(D)] the \emph{direct law} at each square, with a scalar \(c_{i,n}\): \(g_n \lambda_{i+1,n}\delta_{i,n} = (g_{n+2}/\bar c_{i,n})\, \delta_{i,n+1}\lambda_{i,n+1}\) on \(H_{i,n+1}\), where a quotient by \(0\) means \(0\);
\item[(A)] the \emph{adjoint law} at each square, with a scalar \(c^*_{i,n}\): every \(v \in H_{i,n+2}\) has an \(x \in H_{i,n+1}\) with \(\langle \lambda_{i,n+1}y, v\rangle = \langle y, x\rangle\) for all \(y\), and \(\langle \delta_{i,n}x, y'\rangle = c^*_{i,n} \langle \delta_{i,n+1}v, \lambda_{i+1,n}y'\rangle\) for all \(y' \in H_{i+1,n}\).
\end{itemize}

\begin{problem}\label{pr:toy-square}
Let \(m = 2\), every \(H_{i,n} = \Hc = \C\) with the usual inner product and identity embeddings, \((g_0, g_1, g_2) = (2, 3, 5)\), and let the maps be multiplication by \(\lambda_{0,0} = 2\), \(\lambda_{1,0} = 3\), \(\lambda_{0,1} = 4\), \(\delta_{0,0} = 2\), \(\delta_{1,0} = 3\), \(\delta_{0,1} = 1\).  Then (F) holds; the only square is \((0,0)\), its direct law holds exactly with \(c_{0,0} = 5/3\), and its adjoint law exactly with \(c^*_{0,0} = 8/3\).
\end{problem}

\begin{problem}\label{pr:toy-gram}
In Problem~\ref{pr:toy-square}, \(\lambda_0^{(2)} = 48\), and
\(48^2 = 2304 = c_{0,0}c^*_{0,0}\,(g_0g_1^2/g_2)\,(g_0\lambda_{0,0}^2)(g_0\lambda_{1,0}^2)\).
\end{problem}

In Problems~\ref{pr:row-commute}--\ref{pr:shorten} the triangle is arbitrary and satisfies (F), (D), (A).  For \(i + n < m\) put \(c_i^{[n]} := c_{i,0}c_{i,1}\cdots c_{i,n-1}\) and \(c_i^{*[n]} := c^*_{i,0}\cdots c^*_{i,n-1}\) (empty products are \(1\)); Problems~\ref{pr:row-commute}, \ref{pr:row-adjoint} and~\ref{pr:shorten} assume \(i + n < m\).

\begin{problem}\label{pr:row-commute}
\(\lambda_{i+1}^{(n)}\delta_{i,0}\lambda_{i,0} = \dfrac{g_{n+1}}{g_1g_0\,\overline{c_i^{[n]}}}\,\delta_{i,n}\lambda_i^{(n+1)}\) on \(H_{i,0}\).
\end{problem}

\begin{problem}\label{pr:row-injective}
For \(i + n \le m\), \(\lambda_i^{(n)}\) is injective on \(\Hc\).  If \(\Hc \ne 0\) and \(i + n < m\), then \(c_i^{[n]} \ne 0\).
\end{problem}

\begin{problem}\label{pr:row-adjoint}
For \(x \in \Hc\) and \(v \in H_{i,n+1}\): \(\langle v, \lambda_{i,n}\lambda_i^{(n)}x\rangle = c_i^{*[n]}\,\langle \delta_{i,n}v, \lambda_{i+1}^{(n)}x\rangle\).
\end{problem}

\begin{problem}\label{pr:shorten}
For \(x, y \in \Hc\) and \(x' := \delta_{i,0}\lambda_{i,0}x \in \Hc\):
\(\langle \lambda_i^{(n+1)}x, \lambda_i^{(n+1)}y\rangle = c_i^{[n]}c_i^{*[n]}\,\dfrac{g_0g_1g_n}{g_{n+1}}\,\langle \lambda_{i+1}^{(n)}x', \lambda_{i+1}^{(n)}y\rangle\).
\end{problem}

In the chain, \(H_{i,n}\) consists of forms on a Hermitian space of rank \(r + nt\) whose parameter has \(n\) blocks of length \(t\), labelled \(i+1, \dots, i+n\).  In row \(0\) these spaces are \(Y_0, \dots, Y_m\); the first space of row \(i\) is \(Y_0\) with its form multiplied by \((-1)^i\).  On labels, the lift \(\lambda_{i,n}\) adds block \(i+n+1\) and the descent \(\delta_{i,n}\) removes block \(i+1\).  So the square \((i,n)\) removes block \(i+1\) and adds block \(i+n+2\); we say it carries the pair \((i+1, i+n+2)\).

\begin{problem}\label{pr:pairs}
As \((i,n)\) runs over the squares of a triangle of size \(m\), the pair \((i+1, i+n+2)\) runs over the pairs \(1 \le i' < j' \le m\), each once; there are \(m(m-1)/2\) squares.  For \(m = 3\) the squares \((0,0), (0,1), (1,0)\) carry \((1,2), (1,3), (2,3)\).
\end{problem}

\section{The triangle of the chain and its scalars}\label{sec:scalars}

For an \(m\)-fold twist \(\vartheta\) let \(\mathfrak m := (\pi_M)\prod_{w \in Q_\vartheta} w\), so \(S_{\mathfrak m} = \cM \cup Q_\vartheta\), and in this sheet \(\mathfrak R := \Oc_E/\mathfrak m\), both depending on \(\vartheta\); \(\mathfrak R \cong \prod_{w \in \cM \cup Q_\vartheta} k_w\), as \(p_0, p_1\) split completely in \(E\) and \(\Nm Q_\vartheta\) is prime to \(M\).  In the triangle of the chain, \(\Hc = H(Y_0)\), the node \(H_{i,n}\) is a space of forms on a space of rank \(r + nt\), row \(0\) runs through \(Y_0, \dots, Y_m\), the lift \(\lambda_{i,n}\) is a theta lift built from the character \(\mu^0_{i+n+1}\vartheta^{i+n+1}\) (times a power of the splitting character \(\upsilon\), depending on the ranks) with the indicator of \(\binom{1}{0}\) modulo \(\mathfrak m\) as finite datum, and \(g_n(\vartheta) := \#\GL_{r+nt}(\mathfrak R)\).  The descents are transposed theta lifts; in each row the proof chooses their polynomials at infinity and the finite datum of \(\delta_{i,0}\) so that \(\delta_{i,0}\) is adjoint to \(\lambda_{i,0}\) on \(\Hc\).  For a function \(p\) on matrices \(\eta\) of a fixed shape over \(\mathfrak R\), the \emph{next datum} is the function on the matrices with one more block of rows and of columns, of size \(t\), that equals \(p(-\eta)\), divided by the number of possible blocks \(\eta'\), at \(\bigl(\begin{smallmatrix}\eta & \eta'\\ 0 & 1\end{smallmatrix}\bigr)\), and \(0\) at all other matrices; the finite datum of \(\delta_{i,n+1}\) is the next datum of that of \(\delta_{i,n}\).

\emph{Realization of the triangle.}  The proof shows (F), and (D), (A) at each square with functions \(c_{i,n}\), \(c^*_{i,n}\) of the \(m\)-fold twist, as follows.  Theorem~\ref{c3-thm:exchange}, on the sum of the two middle Hermitian spaces, gives a finite square law for the two composites around the square on point data; a finite summation turns it into (D) for a datum and its next datum, with \(\lvert c_{i,n}(\vartheta)\rvert^2 = C_{i,n}\,\ex_t(\cM \cup Q_\vartheta, \nu_{i,n}(\vartheta))\), \(C_{i,n} \in \R\) independent of \(\vartheta\) and \(\nu_{i,n}(\vartheta)\) as in Problem~\ref{pr:square-char}.  The proof also provides the adjoint vectors in (A), with \(c^*_{i,n} \doteq 1\) in the sense of Definition~\ref{def:doteq}; one of its factors is the scalar of the reversed square, which combines an exchange with partners of equal rank with the inverse of the Fourier constant of Problem~\ref{pr:gauss}.  \emph{Status}: read from the proof.  For joint harmonics see \cite{Li90, Ho89}, for the seesaw principle \cite{Ku84}.

\begin{definition}\label{def:doteq}
For functions \(c\) from \(m\)-fold twists to \(\C\) and \(f\) to \(\R\), \(c \doteq f\) means: there are \(C' \in \R\) (possibly \(0\)) and \(e' \in \Z\) with \(\lvert c(\vartheta)\rvert^2 = C'(\Nm Q_\vartheta)^{e'} f(\vartheta)^2\) for every \(m\)-fold twist \(\vartheta\).
\end{definition}

The relation \(c \doteq f\) allows \(c\) to vanish identically; if \(c\) is nonzero at one twist, its zeros are exactly those of \(f\).

\begin{definition}\label{def:rho}
For \(1 \le i < j \le m\), \(\nu := \mu_j^0/\mu_i^0\) and an \(m\)-fold twist \(\vartheta\),
\[
  \rho_{ij}(\vartheta) := \prod_{z=1}^{t} \frac{\lvert L^{\cM \cup Q_\vartheta}\bigl(z, (\nu\vartheta^{j-i})^{-1}\bigr)\rvert}{\lvert L^{\cM}(z, \nu^{-1})\rvert} ,
\]
a quotient by \(0\) being \(0\).
\end{definition}

The values \(L^S(z, \nu)\) in Definition~\ref{def:rho} and in the exchange factors below are finite and nonzero, by the nonvanishing stated in Chapter~\ref{c3-chap}, Section~\ref{c3-sec:integrals}: the characters have weight \(\pm(j-i)\Delta \ne 0\) and a tame part (Problems~\ref{pr:square-char} and~\ref{pr:rho-pair}).

\begin{problem}\label{pr:square-ranks}
In the square \((i,n)\) the Hermitian spaces underlying the middle corners \(H_{i,n+1}\), \(H_{i+1,n+1}\) have rank \(r + (n+1)t\) each, and those underlying the partners \(H_{i,n+2}\), \(H_{i+1,n}\) have ranks \(r + (n+2)t\) and \(r + nt\).  So in Theorem~\ref{c3-thm:exchange}, with \(V^0\) the orthogonal sum of the two middle Hermitian spaces, \(n_0 = 2(r + (n+1)t)\) and \(h = t\).  In the reversed square the partners are the spaces underlying \(H_{i,n+1}\), \(H_{i+1,n+1}\), of equal rank, so \(h = 0\) and \(\ex_0 = 1\).
\end{problem}

For the data, \(8m \le 16\sqrt x < x^{11} \le T(s) \le q_1 d\) (Definition~\ref{da-def:data}, Problem~\ref{da-pr:T}) and \(q_1 d \mid q_2\) (Chapter~\ref{da-chap}).

\begin{problem}\label{pr:square-char}
Let \(\nu_{i,n}(\vartheta) := \bigl((\mu^0_{i+n+2}/\mu^0_{i+1})\vartheta^{n+1}\bigr)^{-1}\).  It has weight \((n+1)\Delta\); moreover \((n+1)\Delta \ge r + (n+1)t\) and \(n_0 \le 2(n+1)\Delta\).  Its tame part is \(\omega_w^{-(n+1)}\) at \(w \in \cM\), \(\psi_w^{-(n+1)}\) at \(w \in Q_\vartheta\), trivial elsewhere, and nontrivial at every \(w \in \cM \cup Q_\vartheta\).
\end{problem}

\begin{problem}\label{pr:square-scalars}
If \(c \doteq f\) and \(c' \doteq f'\), then \(cc' \doteq ff'\).  For every square, \(c_{i,n} \doteq \rho_{i+1,\,i+n+2}\), with exponent \(e' = t^2\), and \(c^*_{i,n} \doteq 1\).
\end{problem}

The powers of \(\#\mathfrak R\) in the adjoint law come from finite Fourier transforms \(\hat F(\mathbf v') := \sum_{\mathbf v} F(\mathbf v)\,\mathbf e(\operatorname{tr}(\mathbf v'\mathbf v))\), with \(\mathbf e\) a primitive additive character of \(\mathfrak R\) (nontrivial on each field factor).  For \(\mathfrak R = \prod_w k_w\), with primitive characters \(\mathbf e\) (additive) and \(\chi\) (multiplicative, \(0\) on nonunits: every \(a \ne 0\) has some \(\chi(1+ca) \notin \{0,1\}\)) there is \(G \in \C\) with \(\lvert G\rvert^2 = \#\mathfrak R^{t^2}\) such that every function \(F\) on \(t \times t\) matrices with \(F(k\mathbf v) = \chi(\det k)F(\mathbf v)\) for invertible \(k\) has \(\hat F(\mathbf v') = G\,\chi^{-1}(\det \mathbf v')\,F(1)\), with \(\chi^{-1}\) also extended by \(0\).  \emph{Status}: for the Gauss sum of a determinant character over a finite field see \cite{Ko63}; for \(\mathfrak R = \prod_w k_w\) such an \(F\) vanishes off the invertible matrices and is a product over \(w\), so the identity follows from the case of fields.  At a node of rank \(r+nt\), let \(\iota=\binom{1}{0}\) be the inclusion into rank \(r+(n+1)t\).  The proof compares two sums over matrices \(x\), each against \(L((\iota\;x))\), where \((\iota\;x)\) denotes adjoining the columns of \(x\) to \(\iota\) and \(L(k\mathbf v)=\chi(\det k)L(\mathbf v)\) for invertible \(k\).  The sum weighted by the Fourier transform of the conjugate of the next datum of \(p\) is \(G\,\#\mathfrak R^{(r+nt)t}\) times the sum weighted by the next datum of the Fourier transform of \(\bar p\): \(G\) comes from a corner block of size \(t\), and \(\#\mathfrak R^{(r+nt)t}\) from the sum over a free block of size \((r+nt)\times t\).  The reversed square uses the inverse of \(G\,\#\mathfrak R^{(r+nt)t}\) (in the proof of (A)).

\begin{problem}\label{pr:next}
Take the added block of size \(1\).  Over \(\mathfrak R = \F_3\), with \(p\) the indicator of \(\{1\}\) on \(1 \times 1\) matrices, the next datum is \(1/3\) at the three matrices \(\bigl(\begin{smallmatrix}-1 & \eta'\\ 0 & 1\end{smallmatrix}\bigr)\) (\(\eta' \in \F_3\)) and \(0\) elsewhere.  With \(\mathbf e(x) = e^{2\pi i x/3}\), its Fourier transform at \(\bigl(\begin{smallmatrix}A & B\\ C & D\end{smallmatrix}\bigr)\) is \(1_{C=0}\,\mathbf e(-A+D)\).
\end{problem}

\begin{problem}\label{pr:gauss}
For \(\mathbf e\) as in Problem~\ref{pr:next} and the nontrivial character \(\chi\) of \(\F_3^\times\), the Gauss sum \(\sum_{x \in \F_3^\times}\chi(x)\,\mathbf e(x)\) is \(i\sqrt3\).  In general \(\lvert G\cdot\#\mathfrak R^{(r+nt)t}\rvert^2 = \#\mathfrak R^{t(t + 2(r+nt))}\); formally, for \(\mathfrak R = \F_3\) and \(r + nt = t = 1\), this is \(\lvert 3i\sqrt3\rvert^2 = 27\).
\end{problem}

\begin{problem}\label{pr:flag}
Over \(\F_q\), the invertible block upper triangular matrices with diagonal blocks of sizes \(r, t, \dots, t\) (\(j\) blocks \(t\)) form a group of order \(\#\GL_r(\F_q)\,\#\GL_t(\F_q)^j\,q^{jrt + t^2 j(j-1)/2}\).  For blocks \((1,2)\) over \(\F_2\) the order is \(24\); for \(r = 1\), \(t = 5\) the exponent is \(5, 35, 90\) for \(j = 1, 2, 3\).
\end{problem}

\section{The Gram identity of the chain}\label{sec:gram}

\begin{definition}\label{def:rallis-product}
For an \(m\)-fold twist \(\vartheta\), let \(\Rc_\vartheta\) be the endomorphism of \(T\)
\[
  \Rc_\vartheta := \Rc_1(\vartheta; Y_0)\,\Rc_2(\vartheta^2; Y_0) \cdots \Rc_m(\vartheta^m; Y_0) .
\]
\end{definition}

Each factor preserves \(T\) (Corollary~\ref{c2-cor:equal}).  In this sheet \(\cmat(\vartheta) := \#\mathfrak R^{r(r+t)}\), the number of \((r+t) \times r\) matrices over \(\mathfrak R\).  The inner product formula on the first lift of each row (Problem~\ref{pr:range}), with the point datum, gives: for \(i < m\) there is \(C_i \in \C\), independent of \(\vartheta\), with
\[
  g_0\,\langle \lambda_{i,0}x, \lambda_{i,0}y\rangle = \frac{C_i}{\cmat(\vartheta)}\,\langle x, \Rc_{i+1}(\vartheta^{i+1}; Y_0)\,y\rangle \qquad (x \in H(Y_0),\ y \in T).
\]
\emph{Status}: read from the proof.

\begin{problem}\label{pr:full-row}
In a triangle of size \(m\) with (F), (D), (A), let \(T' \subseteq \Hc\) be a subspace, \(\mathcal E_i\) endomorphisms of \(T'\) and \(C_i\) scalars with \(g_0\langle \lambda_{i,0}x, \lambda_{i,0}y\rangle = C_i\langle x, \mathcal E_i y\rangle\) for \(i < m\), \(x \in \Hc\), \(y \in T'\).  Then for \(i + n = m\), \(x \in \Hc\) and \(y \in T'\)
\[
  \langle \lambda_i^{(n)}x, \lambda_i^{(n)}y\rangle = c_{(i)}\,\frac{g_0g_1^n}{g_n}\,\langle x, \mathcal E_i\mathcal E_{i+1}\cdots\mathcal E_{m-1}\,y\rangle ,
\]
where \(c_{(i)}\) is the product of \(C_i, \dots, C_{m-1}\) and of all \(c_{i',n'}c^*_{i',n'}\) with \(i' \ge i\).
\end{problem}

\begin{problem}\label{pr:chain-row}
For the triangle of the chain there is \(c_\triangle \doteq \prod_{1 \le i < j \le m}\rho_{ij}\) with
\[
  \langle \lambda_0^{(m)}x, \lambda_0^{(m)}y\rangle = c_\triangle(\vartheta)\,\frac{g_0g_1^m}{g_m\,\cmat^m}\,\langle x, \Rc_\vartheta y\rangle \qquad (x \in H(Y_0),\ y \in T).
\]
\end{problem}

On \(H(Y_0)\) we write \(\Theta_\vartheta\) for the lower chain with values in algebraic forms on \(Y_m\), and on \(T\) for its restriction, as in Chapter~\ref{ob-chap}.  Its target consists of algebraic forms of the prescribed weight on \(Y_m\) on which the final level \(\Klast(\vartheta) \subseteq K^\circ(Y_m)\) acts by its character (Definition~\ref{ob-def:final-level}): the elements of \(K^\circ(Y_m)\) that reduce to \(1\) at each \(w \in \cM\) and to a block upper triangular matrix, with blocks \(r, t, \dots, t\), at each \(w \in Q_\vartheta\).  The proof shows
\(
  [K^\circ(Y_m) : \Klast(\vartheta)] = \prod_{w \in \cM}\#\GL_{r+mt}(k_w)\,\prod_{w \in Q_\vartheta}[\GL_{r+mt}(k_w) : \mathrm{Par}_w]
\),
with \(\mathrm{Par}_w\) the group of Problem~\ref{pr:flag} over \(k_w\), and that for \(x, y \in H(Y_0)\)
\[
  \langle \Theta_\vartheta x, \Theta_\vartheta y\rangle = \cwt\,[K^\circ(Y_m) : \Klast(\vartheta)]\,\langle \lambda_0^{(m)}x, \lambda_0^{(m)}y\rangle ,
\]
with \(\cwt\) a positive integer independent of \(\vartheta\).  \emph{Status}: read from the proof.

\begin{problem}\label{pr:theta-injective}
\(\Theta_\vartheta\) is injective on \(H(Y_0)\), in particular on \(T\).
\end{problem}

\begin{problem}\label{pr:last-index}
For every \(m\)-fold twist \(\vartheta\),
\[
  [K^\circ(Y_m) : \Klast(\vartheta)]\,\frac{g_0g_1^m}{g_m} = \bigl(\cmat(\vartheta)\,c_{Q_\vartheta}\bigr)^m\,(\Nm Q_\vartheta)^{-mrt - t^2m(m-1)/2}
\]
times a positive number independent of \(\vartheta\).
\end{problem}

\begin{problem}\label{pr:assemble}
There is \(c\) with \(c \doteq \prod_{1 \le i < j \le m}\rho_{ij}\) and \(\langle \Theta_\vartheta x, \Theta_\vartheta y\rangle = c(\vartheta)\,c_{Q_\vartheta}^m\,\langle x, \Rc_\vartheta y\rangle\) for all \(m\)-fold twists \(\vartheta\) and \(x, y \in T\).
\end{problem}

Problems~\ref{pr:row-commute}--\ref{pr:pairs}, \ref{pr:square-scalars} and~\ref{pr:full-row}--\ref{pr:assemble}, with the realization of the triangle, the first-edge identity and the description of the last level, prove the following.

\begin{theorem}\label{thm:chain-gram}
There is a function \(c\) from \(m\)-fold twists to \(\C\) with \(c \doteq \prod_{1 \le i < j \le m}\rho_{ij}\) such that for every \(m\)-fold twist \(\vartheta\) and all \(x, y \in T\),
\[
  \langle \Theta_\vartheta x, \Theta_\vartheta y\rangle = c(\vartheta)\,c_{Q_\vartheta}^m\,\langle x, \Rc_\vartheta y\rangle .
\]
\end{theorem}

The data satisfy \(t \le \Delta\) and \(m\Delta \le 4b\) (as \(m\Delta \le 4\sqrt x\,(3x^5)^8 < x_1 < b\) by Definition~\ref{da-def:data}).

\begin{problem}\label{pr:rho-pair}
\(\rho_{ij}(1) = 1\).  The character \(\mu_j^0/\mu_i^0\) has weight \(-(j-i)\Delta\), with \(t \le (j-i)\Delta \le 4b\), and tame part \((\omega_w^{j-i})\), nontrivial exactly at the \(w \in \cM\); so each \(P_z := (\mu_j^0/\mu_i^0, z)\), \(1 \le z \le t\), is a ramified critical pair (Definition~\ref{ob-def:pair}), and \(\rho_{ij}(\vartheta) = \prod_{z=1}^t R_{P_z}(\vartheta^{j-i}, 1)\).
\end{problem}

Problem~\ref{pr:rho-pair} identifies the unprimed ratios of Problem~\ref{dch-pr:count-ratios}; Theorem~\ref{dch-thm:small-ratios}, used in Chapter~\ref{tc-chap}, bounds for a family with small ratios the height of the inverses of \((\Nm Q_{\theta_\iota})^{e/2}\prod_{i<j}\rho_{ij}(\theta_\iota)\), every \(e \in \Z\).  We have proved Theorem~\ref{thm:chain-gram} from three statements read from the proof; the exchange identity enters through the square laws.  Chapter~\ref{dch-chap} uses it to factor \(Z(\vartheta)/Z(1)\).

\paragraph{In the Lean formalization.}
\begin{sloppypar}
In the Lean proof \cite{GL} (\tl{grh/\allowbreak src/\allowbreak New/\allowbreak Determinant/\allowbreak Chain/}), Theorem~\ref{thm:chain-gram} is \tl{New.\allowbreak Data.\allowbreak exists\_\allowbreak inner\_\allowbreak chain\_\allowbreak eq}; \(\rho_{ij}\), \(\Rc_\vartheta\), \(\doteq\) are \tl{quotientRatio}, \tl{rallisProduct}, \tl{HasAbsUpToNormPower} (for the indices used here).  The realization is \tl{exists\_\allowbreak isTriangle} (flag \tl{false}), using \tl{exists\_\allowbreak triangleLift\_\allowbreak triangleTransposedLift\_\allowbreak comm} for (D) and \tl{exists\_\allowbreak triangleTransposedLift\_\allowbreak adjoint\_\allowbreak comm} for the identity in (A).  The first-edge identity is a step in the proof of \tl{exists\_\allowbreak inner\_\allowbreak weightedChain\_\allowbreak eq}.  The comparison at the last level is \tl{inner\_\allowbreak lowerChainAtLevel}, and the product formula for the index of \(\Klast(\vartheta)\) is \tl{stdLevelIndex\_\allowbreak lastLevel\_\allowbreak eq\_\allowbreak prod}, both in \tl{ChainNorms.lean}.
\end{sloppypar}

\section*{Hints}
\begin{description}
\item[\ref{pr:row-injective}] For injectivity, induction on \(n\), passing a kernel vector to row \(i+1\) by Problem~\ref{pr:row-commute} and (F).  For \(c_i^{[n]} \ne 0\), take a nonzero vector of \(\Hc\); a quotient by \(0\) is \(0\).
\item[\ref{pr:row-adjoint}] Induction on \(n\); for the step use (A).
\item[\ref{pr:shorten}] Problem~\ref{pr:row-adjoint} with \(v = \lambda_i^{(n+1)}x\), then Problem~\ref{pr:row-commute} for \(\lambda_{i+1}^{(n)}x'\).
\item[\ref{pr:square-char}] The weights \(\kappa_j\) of Chapter~\ref{ob-chap} increase by \(\Delta\).  For the rank bound use \(r \ge 3\), \(t = 5r\), \(\Delta = 2r^8\); for nontriviality use \(0 < n+1 \le m < q_2\) and the \(m\)-fold condition.
\item[\ref{pr:square-scalars}] Problem~\ref{c3-pr:ex-ratio} with \(S = \cM\), \(Q = Q_\vartheta\), \(h = t\); the untwisted value \(\ex_t(\cM, \nu_{i,n}(1))\) is a constant.  Note \(\#\mathfrak R = \Nm Q_\vartheta \prod_{w \in \cM} q_w\).
\item[\ref{pr:full-row}] Induction on \(n\) from row \(i+1\) to row \(i\), with Problem~\ref{pr:shorten}, and (F) to apply the first-edge relation to \(x\) and \(\mathcal E_{i+1}\cdots\mathcal E_{m-1}y\).
\item[\ref{pr:chain-row}] Problem~\ref{pr:full-row} with \(\mathcal E_i = \Rc_{i+1}(\vartheta^{i+1}; Y_0)/\cmat\), and Problems~\ref{pr:pairs} and~\ref{pr:square-scalars}.
\item[\ref{pr:last-index}] Factor over \(w \in \cM \cup Q_\vartheta\); at \(w \in Q_\vartheta\) use Problem~\ref{pr:flag}, the stabilizer of an injective \((r+t) \times r\) matrix, and Problem~\ref{c1-pr:left-inverse}.
\item[\ref{pr:assemble}] Problems~\ref{pr:chain-row} and~\ref{pr:last-index}; for \(A \in \C\) and \(k \in \Z\), \(\vartheta \mapsto A(\Nm Q_\vartheta)^k\) satisfies \(\doteq 1\).
\end{description}

\begin{thebibliography}{Li90}
\bibitem[GL]{GL} L.~Sernau, \emph{GRH for Dirichlet \(L\)-functions, in Lean~4}, a shorter rewriting of the Lean proof, text by Claude, \url{https://github.com/lsernau/grh-lean}, commit \texttt{72fead1}, readable edition \texttt{grh/}.
\bibitem[Ho89]{Ho89} R.~Howe, \emph{Transcending classical invariant theory}, J.~Amer. Math. Soc. \textbf{2} (1989), 535--552.
\bibitem[Ko63]{Ko63} T.~Kondo, \emph{On Gaussian sums attached to the general linear groups over finite fields}, J.~Math. Soc. Japan \textbf{15} (1963), 244--255.
\bibitem[Ku84]{Ku84} S.~S.~Kudla, \emph{Seesaw dual reductive pairs}, in: Automorphic Forms of Several Variables (Katata, 1983), Progr. Math. \textbf{46}, Birkh\"auser, 1984, 244--268.
\bibitem[Li90]{Li90} J.-S.~Li, \emph{Theta lifting for unitary representations with nonzero cohomology}, Duke Math. J. \textbf{61} (1990), 913--937.
\end{thebibliography}

\end{document}
